\documentclass[journal]{IEEEtran}
\usepackage{graphicx,booktabs,amsmath,amssymb,url,xcolor,microtype}
\usepackage[hidelinks]{hyperref}
\graphicspath{{../assets/figures/}}

\title{Jev-Grounded Adaptive Temporal Granularity for Causal Video Generation:\\
An Image-Grounded Off-Device Oracle with Recovery, Recalibration, and an Independent Endpoint}

\author{Hossen Md Jisan\\
Department of Artificial Intelligence, NUIST\\
Email: jisanahammed1952@gmail.com}

\begin{document}
\maketitle

\begin{abstract}
Adaptive temporal granularity---deciding when a causal video generator should produce one frame at a
time (fine) versus several frames per chunk (coarse)---requires three preconditions: a real
mode-dependent cost gap, a pre-generation signal, and an endpoint independent of the compared modes.
We audit all three with an off-device, image-grounded System One oracle (Jev~1.14). We contribute
(i)~a matched break-even law over hardware and operating point, (ii)~a paired test of image-grounded
versus text-grounded routing signals, (iii)~a validated independent image-grounded endpoint, and
(iv)~two honest negatives: an image signal that does not beat text, and a stability-gated
recalibration policy that does not help. On a matched factorial (same prompts, same modes, T4 and
A100) we show the adaptive outcome flips with the coarse operating point alone, explaining a
confound in prior work. The image-grounded oracle is a poor \emph{signal} but an excellent
\emph{endpoint} (Spearman $\rho{=}{+}1.000$ against CLIP subject consistency).
\end{abstract}

\begin{IEEEkeywords}
adaptive granularity, temporal granularity, causal video generation, video world models,
autoregressive diffusion, off-device oracle, System One, efficient inference.
\end{IEEEkeywords}

\section{Introduction}
Causal (autoregressive) video generators can trade compute for quality by changing their temporal
granularity: frame-wise generation is accurate but expensive, while chunk-wise generation is cheaper
on some hardware but drifts. Prior work adapts \emph{denoising steps} with model-internal signals.
We ask whether an \emph{off-device}, image-grounded System One model can drive the granularity
decision and certify its outcome.

This paper is a \emph{characterization with one strong positive}: we find the image oracle is a poor
routing signal but an excellent independent endpoint, and we isolate a hardware$\times$operating-point
break-even law that resolves a confound in the earlier audit of this problem~\cite{jisan2026audit}.

\section{Preconditions}
\textbf{P1 (Cost).} Coarse must be cheaper than fine \emph{on the deployment hardware and at the
chosen operating point}. \textbf{P2 (Signal).} A pre-generation signal must identify coarse-unsafe
segments. \textbf{P3 (Endpoint).} Quality must be certified by an endpoint independent of the
compared modes. \textbf{P4 (Calibration).} The oracle's confidence must be calibrated enough to gate
on.

\section{Method}
We use Jev~1.14 as an off-device oracle. Text states pass a caption and motion statistics; image
states pass the last three context frames as base64 assets with the header
\texttt{TypeSafe-Danger-Zone-Images:~true}. The backbone is Causal-rCM Wan2.1-T2V-1.3B (commit
\texttt{ed3cb14}). Modes: fine \texttt{c1-1} ($S$ steps), cheap coarse \texttt{c3-3 s1}, expensive
coarse \texttt{c3-3 s4}, and hold. Coarse is cheaper iff
$\;s_{\text{nfe}}^{\text{coarse}}/s_{\text{nfe}}^{\text{fine}} < \mathrm{NFE}_{\text{fine}}/\mathrm{NFE}_{\text{coarse}}$.

\section{Results}

\subsection{E1 --- Break-even is hardware- and operating-point-conditional}
A matched matrix over $\{T4,A100\}\times\{$fine s1/s2/s4$\}\times\{$coarse s1/s2/s4$\}$
(Fig.~\ref{fig:e1}) shows: on T4 the cheap coarse \texttt{c3-3 s1} beats fine only at fine
$s2$/$s4$, and the expensive \texttt{c3-3 s4} is never cheaper; on A100 the cheap coarse beats all
fine modes. This resolves the prior audit's confound, where T4 was tested with the \emph{expensive}
coarse and A100 with the \emph{cheap} coarse.

\subsection{E2 --- Image-grounded Jev does not beat text Jev as a signal}
On the same 80 segments, image and hybrid states are \emph{worse} than the text state
(paired $\Delta$AUROC: difficulty $-0.110$, drift $-0.108$, motion $-0.195$ with a 95\% CI excluding
zero). Text-1.13 $\approx$ text-1.14, so the deficit is the modality, not the model version
(Fig.~\ref{fig:e2}). This is an honest negative.

\subsection{E3 --- The adaptive outcome flips with the operating point alone}
Same A100, same 12 prompts, only the coarse step differs. With the expensive coarse the online
policy is forced to 92\% fine; with the cheap coarse it uses 21\% fine and beats \texttt{fixed\_fine}.

\subsection{E4 --- Recovery$+$hysteresis and the switch seam}
Recovery$+$hysteresis reduces switches on both GPUs without worsening regret, and a switch
\emph{seam} exists (switch-adjacent L1 $>$ interior L1, ratio $>1$).

\subsection{E5 --- Image-grounded Jev is a validated independent endpoint}
Across 6 policies and 48 videos, the image-grounded endpoint agrees with CLIP subject consistency at
Spearman $\rho{=}{+}1.000$ (same-subject) and $+0.943$ (consistency) (Fig.~\ref{fig:e5}).

\subsection{E6 --- Matched T4 boundary}
With fine \texttt{c1-1 s1} the cheap coarse saves $-4.57$~s on T4, so adaptive routing cannot save
time regardless of the signal: the T4 negative is an operating-point boundary, not a signal failure.

\subsection{E7 --- Confidence-gating is neutral}
Gating on Jev's confidence does not improve the regret--cost frontier ($|\Delta|<0.005$ at all budgets).

\subsection{E8 --- Stability-gated recalibration does not help}
Injecting a high-step fine segment on drift is worse on both GPUs: A100 recalibrate minus full
$=$ refL1 $+0.0112$, DiT $+2.04$~s; T4 recalibrate minus adaptive\_warp $=$ refL1 $+0.0202$,
DiT $+12.73$~s (Fig.~\ref{fig:e8}). An honest negative.

\subsection{E9 --- Claim-to-artifact audit}
Every headline number is re-derived from stored artifacts: 18/18 PASS.

\section{Discussion and Limitations}
The image-grounded oracle is a poor \emph{signal} but an excellent \emph{endpoint}; the break-even
law is the reusable tool. RQ2 (the cost--quality trade-off) is hardware- and operating-point-
conditional and is our only headline negative besides E2/E8, which we report openly. Jev~1.14 images
are ``Dreamer-ready''; we pin the version. The signal test uses 80 paired segments.

\section{Conclusion}
An off-device, image-grounded System One model can certify adaptive temporal granularity even when it
cannot drive it, and the cost side obeys a matched break-even law over hardware and operating point.

\begin{figure}[t]\centering\includegraphics[width=\linewidth]{fig_e1_breakeven.png}
\caption{E1 break-even over hardware and operating point.}\label{fig:e1}\end{figure}
\begin{figure}[t]\centering\includegraphics[width=\linewidth]{fig_e2_signal.png}
\caption{E2 image/hybrid Jev do not beat text Jev as a signal.}\label{fig:e2}\end{figure}
\begin{figure}[t]\centering\includegraphics[width=\linewidth]{fig_e5_endpoint.png}
\caption{E5 image-Jev endpoint agrees with CLIP.}\label{fig:e5}\end{figure}
\begin{figure}[t]\centering\includegraphics[width=\linewidth]{fig_e8_recalibration.png}
\caption{E8 recalibration is not beneficial on either GPU.}\label{fig:e8}\end{figure}

\bibliographystyle{IEEEtran}
\bibliography{references}
\end{document}
